\documentclass[11pt,a4paper]{article}

\usepackage[margin=2.4cm]{geometry}
\usepackage{amsmath,amssymb,amsthm}
\usepackage{booktabs}
\usepackage{array}
\usepackage{enumitem}
\usepackage{xcolor}
\usepackage{tikz}
\usepackage{listings}
\usepackage[hidelinks]{hyperref}
\usepackage{caption}

\usetikzlibrary{arrows.meta,positioning,fit,backgrounds}

\definecolor{accent}{HTML}{1D9E75}
\definecolor{lightbox}{HTML}{EEEDFE}
\definecolor{codebg}{HTML}{F5F5F2}

\lstset{
  basicstyle=\ttfamily\footnotesize,
  backgroundcolor=\color{codebg},
  frame=single, rulecolor=\color{gray!40},
  breaklines=true, showstringspaces=false,
  keywordstyle=\color{accent}\bfseries,
  commentstyle=\color{gray},
  language=Python, xleftmargin=0pt, framexleftmargin=4pt
}

\newcommand{\R}{\mathbb{R}}
\newcommand{\E}{\mathbb{E}}
\newcommand{\Prob}{\mathbb{P}}
\newcommand{\Sset}{\mathcal{S}}
\newcommand{\Cset}{\mathcal{C}}

\title{\textbf{A Spatio-Temporal Foundation Model for\\ Multi-Pollutant Air Quality Forecasting}\\[0.4em]
\large Design specification, training protocol and evaluation plan}
\author{Design document}
\date{\today}

\begin{document}
\maketitle

\begin{abstract}
\noindent
We specify a spatio-temporal foundation model (AQFM) that consumes irregularly
sampled, block-missing, multi-station air quality observations and emits
calibrated probabilistic forecasts for multiple pollutants simultaneously. The
design is driven by three empirical findings from the Delhi CPCB exploratory
analysis: (i) missingness is \emph{systemic}, with
$\Prob(\text{PM}_{10}\ \text{missing}\mid \text{PM}_{2.5}\ \text{missing})=0.928$,
so gaps are station-wide contiguous blocks rather than isolated dropouts;
(ii) strong cross-pollutant structure exists ($r_{\text{PM}_{2.5},\text{PM}_{10}}=0.84$,
$r_{\text{NO},\text{NO}_x}=0.90$) which can be imposed by construction instead of
learned; and (iii) the 2018--2026 subset achieves $7$--$10\%$ missingness on primary
pollutants, making it the viable pre-training window. We compare four decoder
designs, recommend a query-based unified decoder, and give a complete training
and evaluation protocol.
\end{abstract}

\tableofcontents

%======================================================================
\section{Problem statement and notation}

Let $\Sset=\{1,\dots,S\}$ index monitoring stations and $\Cset=\{1,\dots,C\}$
index channels (pollutants, meteorology, and derived quantities). Observations
arrive on a nominal 15-minute grid but with substantial missingness. Write

\begin{align}
X &\in \R^{S\times C\times T}, &&\text{values (log-transformed, see \S\ref{sec:preproc})}\\
M &\in \{0,1\}^{S\times C\times T}, &&\text{observation indicator, } M_{s,c,t}=1 \text{ iff measured}\\
\tau &\in \R^{T}, &&\text{absolute timestamps in hours}\\
G &\in \R^{S\times 2}, &&\text{station coordinates (lat, lon)}\\
U &\in \R^{S\times D_{\text{st}}}, &&\text{static station covariates.}
\end{align}

Given a context window of length $T_{\text{ctx}}$ and known-future covariates,
the task is to produce the joint predictive distribution

\begin{equation}
p\big(\, X_{s,c,t} \;:\; s\in\Sset,\; c\in\Cset_{\text{tgt}},\; t\in (T_{\text{ctx}}, T_{\text{ctx}}+H] \;\big|\; X_{\cdot,\cdot,\le T_{\text{ctx}}},\, M,\, \tau,\, G,\, U \,\big),
\end{equation}

with $H$ up to 72 hours. We require the model to be well-defined when
(a) arbitrary channels are missing, (b) arbitrary contiguous blocks are missing
across all channels at a station, and (c) the target set $\Cset_{\text{tgt}}$ and
horizon set are chosen at inference time.

\paragraph{Design principle.} Every irregularity is handled by \emph{excluding
tokens}, never by imputing values before the model sees them. Pre-imputation
injects a second model's errors into the training signal and makes the
observation mask unrecoverable.

%======================================================================
\section{Data}

\subsection{Temporal scope}

The EDA is unambiguous: pre-2018 station-years show $80$--$100\%$ missingness
(PM$_{10}$ is $100\%$ missing every year through 2014), while 2018--2026 sits at
$7$--$21\%$ for primary pollutants. We therefore train on \textbf{2018--2026} and
treat pre-2018 as an optional low-weight auxiliary corpus, not as pre-training
data.

\subsection{Scale, and why Delhi alone is not enough}

Delhi's 38 stations yield $\approx 1.44\times 10^{7}$ rows and
$\approx 2.16\times 10^{8}$ scalar observations. Time-series scaling laws differ
from those for language, but the order of magnitude is informative: at the
commonly cited $\sim\!20$ tokens per parameter this supports a model of roughly
$10$--$30$M parameters. The architecture below wants more.

The resolution is the national network: roughly 550 CPCB stations across India
(Figure~\ref{fig:scale} in the accompanying analysis) give 10--15$\times$ the
data. We therefore adopt a two-corpus strategy.

\begin{table}[h]
\centering
\begin{tabular}{@{}lrrl@{}}
\toprule
Corpus & Stations & Approx.\ observations & Role \\
\midrule
All-India CPCB, 2018--2026 & $\sim$550 & $\sim3\times10^{9}$ & Stage-1 pre-training \\
Delhi NCR, 2018--2026      & 38        & $2.2\times10^{8}$   & Stage-2 fine-tuning \\
Delhi NCR, pre-2018        & 38        & --                  & Optional, low weight \\
\bottomrule
\end{tabular}
\caption{Two-corpus training strategy. National pre-training also supplies a
genuine domain-shift test set: Chennai's coastal sea-breeze regime versus
Delhi's continental winter inversion.}
\end{table}

\subsection{Exogenous covariates}

Two additions are likely to matter more than any architectural choice.

\begin{enumerate}[leftmargin=1.4em]
\item \textbf{Boundary layer height.} Beyond roughly 12 hours of lead, forecast
skill is dominated by ventilation rather than by emissions. The CPCB variable
list contains no mixing-depth proxy. Obtain planetary boundary layer height,
10\,m winds, and precipitation from ERA5 reanalysis.
\item \textbf{Known-future meteorology.} Numerical weather prediction output is
available at forecast time. Train with ERA5 as a perfect-prognosis oracle, then
substitute operational NWP at inference; report the skill gap between the two,
since it bounds the achievable operational performance.
\end{enumerate}

The satellite/reanalysis annual-mean PM$_{2.5}$ field supplies static
per-station covariates $U$: climatological burden, urban--rural gradient, and a
spatial prior for stations with short records.

\subsection{Pre-processing}
\label{sec:preproc}

Concentrations are right-skewed and heavy-tailed, so all pollutant channels are
mapped through $x\mapsto\log(1+x)$ before any statistics are computed. Values
are then normalised per (station, channel) using \emph{context-only} statistics:

\begin{equation}
\hat\mu_{s,c}=\frac{\sum_{t\le T_{\text{ctx}}} M_{s,c,t}X_{s,c,t}}{\sum_{t\le T_{\text{ctx}}} M_{s,c,t}},
\qquad
\hat\sigma_{s,c}^2=\frac{\sum_{t\le T_{\text{ctx}}} M_{s,c,t}(X_{s,c,t}-\hat\mu_{s,c})^2}{\sum_{t\le T_{\text{ctx}}} M_{s,c,t}} .
\end{equation}

Statistics computed over the full window would leak the scale of the future into
the input. When $\sum_t M_{s,c,t}$ falls below a threshold the transform falls
back to the identity rather than dividing by a noise estimate.

%======================================================================
\section{Tokenisation}

\subsection{One token per (station, channel, patch)}

Let $P$ be the patch length ($P=16$, i.e.\ 4 hours) and $N=T_{\text{ctx}}/P$.
Define patch slices $x^{(n)}_{s,c}\in\R^{P}$ and $m^{(n)}_{s,c}\in\{0,1\}^{P}$.
The token is

\begin{equation}
\boxed{\;
z_{s,c,n} \;=\; W_{v}\begin{bmatrix} x^{(n)}_{s,c} \\ m^{(n)}_{s,c}\end{bmatrix}
\;+\; e^{\text{ch}}_{c}
\;+\; \phi_{\text{st}}\!\left(\big[G_{s};U_{s}\big]\right)
\;+\; \psi_{\text{t}}\!\left(\bar\tau^{(n)}\right)
\;\in\;\R^{d}\;}
\label{eq:token}
\end{equation}

with $W_v\in\R^{d\times 2P}$ and $\bar\tau^{(n)}$ the mean timestamp of patch $n$.
Three properties follow directly from this construction.

\begin{description}[leftmargin=1.4em,style=nextline]
\item[Channel-count flexibility.] The token carries its own channel identity, so
a station lacking NH$_3$ simply contributes fewer tokens. No padding, no
sentinel values, no per-station bookkeeping.
\item[Station generalisation.] $\phi_{\text{st}}$ is a function of coordinates,
not a lookup table, so a station never seen in training is a valid input.
\item[Mask visibility.] The mask enters the embedding alongside the values, so a
patch that is $30\%$ observed is not confused with a clean one.
\end{description}

\subsection{Three tiers of irregularity}

\begin{table}[h]
\centering
\small
\begin{tabular}{@{}p{4.1cm}p{4.6cm}p{5.6cm}@{}}
\toprule
Irregularity & Mechanism & Rationale \\
\midrule
Partial patch (some steps missing) & Mask channel in Eq.~\eqref{eq:token}, zero-fill values & The embedding learns to discount unobserved positions \\
Fully unobserved patch & Token excluded from every attention key set & Attending to an all-zero token is worse than not attending \\
Non-uniform timestamps & $\psi_{\text{t}}$ is a function of $\tau$, not of position index & A 10-hour gap must look like 10 hours, not one step \\
Asynchronous sources (future) & Continuous-time front end (mTAN, or set attention over $(t,v,c)$ triplets) & Only needed when fusing satellite overpasses or low-cost sensors \\
\bottomrule
\end{tabular}
\caption{Missingness is handled at three levels; the fourth row is a planned
extension, not part of the base model.}
\end{table}

\subsection{Continuous time features}

With log-spaced periods $\{p_k\}_{k=1}^{K_f}$ from 3 hours to one year and
$\omega_k=2\pi/p_k$,

\begin{equation}
\psi_{\text{t}}(\tau) = W_{\tau}\Big[\;
\{\sin \omega_k\tau,\ \cos\omega_k\tau\}_{k=1}^{K_f},\;
\text{cyc}_{24}(\tau),\ \text{cyc}_{168}(\tau),\ \text{cyc}_{8760}(\tau),\;
\log(1+\Delta\tau)\;\Big],
\end{equation}

where $\text{cyc}_{p}(\tau)=(\sin 2\pi\tau/p,\ \cos 2\pi\tau/p)$ and
$\Delta\tau$ is the gap to the previous patch. The explicit gap feature is what
lets the model distinguish ``next patch'' from ``next patch after a ten-hour
outage.''

%======================================================================
\section{Encoder}

\subsection{Axial factorisation}

A single sample contains $S\!\times\!C\!\times\!N$ tokens; at $S=12$, $C=19$,
$N=32$ this is $7{,}296$ tokens, and at the full 38-station Delhi network it
exceeds $23{,}000$. Full self-attention is $O(L^2)$ and out of reach. We
factorise each encoder block into three cheap passes:

\begin{align}
Z &\leftarrow Z + \mathrm{Attn}_{\text{time}}\big(\mathrm{LN}(Z)\big)   && \text{over } n, \text{ batched over } (s,c) \\
Z &\leftarrow Z + \mathrm{Attn}_{\text{var}}\big(\mathrm{LN}(Z)\big)    && \text{over } c, \text{ batched over } (s,n) \\
Z &\leftarrow Z + \mathrm{Attn}_{\text{space}}\big(\mathrm{LN}(Z)\big)  && \text{over } s, \text{ batched over } (c,n) \\
Z &\leftarrow Z + \mathrm{FFN}\big(\mathrm{LN}(Z)\big).
\end{align}

Cost drops from $O\!\big((SCN)^2\big)$ to $O\!\big(SCN(S+C+N)\big)$: at the
figures above, a factor of roughly 300. Stacking $B_{\text{ax}}=4$ such blocks
gives every token an effective receptive field over the whole grid.

\subsection{Masked attention}

Each attention uses a key-validity mask $v_{s,c,n}=\mathbb{1}[\,\sum_t m^{(n)}_{s,c,t}>0\,]$,
implemented as an additive $-\infty$ bias. If every key in a row is invalid the
row is unmasked rather than producing $\mathrm{NaN}$; the mask channel already
tells the model the content is uninformative.

\subsection{Wind-aware spatial bias}

A station 8\,km \emph{upwind} is far more informative than one 8\,km downwind. A
distance-only graph cannot express this asymmetry. Let $d_{ij}$ be the
great-circle distance and $\theta_{ij}$ the bearing from station $i$ to $j$, and
let $\mathrm{WD}_i^{(n)}$, $\mathrm{WS}_i^{(n)}$ be the patch-mean wind direction
and speed at the query station. The additive attention bias is

\begin{equation}
\boxed{\;
b^{(n)}_{ij} = \mathrm{MLP}\Big(
\log(1+d_{ij}),\;
\cos\big(\theta_{ij}-\mathrm{WD}_i^{(n)}\big),\;
\sin\big(\theta_{ij}-\mathrm{WD}_i^{(n)}\big),\;
\mathrm{WS}_i^{(n)},\;
\cos(\cdot)\,\mathrm{WS}_i^{(n)},\;
e^{-d_{ij}/\lambda}
\Big)\;}
\end{equation}

emitted per attention head and added to the spatial attention logits. This is
cheap, and the learned bias is inspectable: plotting which stations the model
attends to under different wind regimes is a publishable result in itself.

\subsection{Per-station latent bottleneck}

Cross-attending $\sim\!1.7$k decoder queries against $\sim\!7$k encoder tokens is
wasteful. Each station's $(C\times N)$ grid is compressed by a Perceiver-style
cross-attention into $L=16$ latents, giving $S\!\cdot\!L$ memory tokens. This
preserves the spatial structure the wind bias operates on while cutting decoder
cross-attention cost by roughly $30\times$.

%======================================================================
\section{Decoder: four designs}
\label{sec:decoders}

This is the central design question. All four options share the encoder above.

\subsection{Design A --- one decoder per pollutant}

Each target pollutant $p$ gets its own decoder $\mathcal{D}_p$ cross-attending to
the shared memory.

\textbf{For.} Pollutants genuinely differ in dynamics. The EDA correlation matrix
makes this concrete: ozone is the only pollutant \emph{positively} correlated
with temperature ($r=+0.34$) while PM$_{2.5}$ sits at $r=-0.49$, reflecting
photochemical production versus accumulation under winter inversion. Separate
decoders also allow adding a pollutant without touching the trunk, and
independent per-pollutant loss weighting.

\textbf{Against.} Decoder parameters scale linearly in $|\Cset_{\text{tgt}}|$; the
substantial structure that \emph{is} shared across pollutants cannot be shared;
and independent heads produce physically incoherent joint output --- nothing
prevents $\widehat{\text{PM}}_{2.5} > \widehat{\text{PM}}_{10}$.

\subsection{Design B --- query-based unified decoder (recommended)}

One decoder. A query is constructed per prediction target:

\begin{equation}
q_{s,c,h} \;=\; e^{\text{tgt}}_{c} \;+\; e^{\text{lead}}_{h} \;+\; \phi_{\text{st}}\!\left([G_s;U_s]\right) \;+\; \psi_{\text{t}}\!\left(\tau^{\text{fut}}_{h}\right),
\end{equation}

followed by $L_{\text{dec}}$ layers of \{query self-attention, cross-attention to
memory, FFN\}. Specialisation lives in the query embedding rather than in
separate weights.

\textbf{For.} Arbitrary target subsets at inference (forecast PM$_{2.5}$ alone
with no wasted compute); arbitrary horizons, including untrained ones if lead
time is encoded continuously; cross-pollutant coherence, because query
self-attention lets the PM$_{2.5}$ prediction attend to the PM$_{10}$
prediction before either is emitted; one weight set to scale when moving from 38
to 550 stations.

\textbf{Against.} Potential under-fitting of a pollutant with genuinely divergent
dynamics --- which is what Design C addresses.

\subsection{Design C --- shared trunk with per-pollutant adapters}

Shared decoder trunk plus FiLM conditioning or low-rank (LoRA-style) adapters per
pollutant. Roughly $2\%$ additional parameters buys most of Design A's
specialisation. \emph{Deploy this only if ablation shows Design B under-fits a
specific pollutant}, most plausibly O$_3$.

\subsection{Design D --- mixture-of-experts decoder}

A router over $k$ experts. Pollutants cluster naturally into families
(particulates, NO$_x$, photochemical, meteorological); MoE lets that clustering be
learned rather than imposed. Worth the complexity only at national scale with
many channels.

\begin{table}[h]
\centering
\small
\begin{tabular}{@{}lcccc@{}}
\toprule
& A: per-pollutant & B: query & C: adapters & D: MoE \\
\midrule
Decoder params vs.\ $|\Cset_{\text{tgt}}|$ & linear & constant & constant$+\epsilon$ & constant \\
Arbitrary target subset at inference & no & yes & yes & yes \\
Arbitrary / unseen horizon & no & yes & yes & yes \\
Joint cross-pollutant coherence & no & yes & yes & yes \\
Per-pollutant specialisation & high & medium & high & high \\
Implementation complexity & low & low & medium & high \\
\midrule
\textbf{Verdict} & baseline & \textbf{recommended} & fallback & at scale \\
\bottomrule
\end{tabular}
\caption{Decoder design comparison. Design B is the recommendation; A is worth
implementing as an ablation baseline because it is the intuitive choice and
should be shown to be dominated.}
\end{table}

%======================================================================
\section{Physical constraints}

Independent heads can emit PM$_{2.5}>$ PM$_{10}$ and $\text{NO}+\text{NO}_2\neq\text{NO}_x$.
Both are free to fix by predicting a base quantity and a bounded fraction:

\begin{equation}
\boxed{\;
\text{PM}_{2.5} = f_{\text{pm}}\cdot \text{PM}_{10},\quad f_{\text{pm}}=\sigma(\tilde f_{\text{pm}})\in(0,1);
\qquad
\text{NO} = f_{\text{nox}}\cdot\text{NO}_x,\quad \text{NO}_2=(1-f_{\text{nox}})\cdot \text{NO}_x. \;}
\end{equation}

The model predicts $\{\log(1+\text{PM}_{10}),\ \tilde f_{\text{pm}},\ \log(1+\text{NO}_x),\ \tilde f_{\text{nox}},\dots\}$;
observables are reconstructed afterwards. Two implementation notes:

\begin{itemize}[leftmargin=1.4em]
\item The logit-fractions are computed in the data pipeline and appended as
extra input channels, so they receive their own normalisation statistics and are
predicted like any other channel.
\item A fraction is treated as observed only when \emph{both} constituents are
observed and the denominator is non-trivial. Otherwise a missing PM$_{10}$
silently manufactures a PM$_{2.5}$ fraction.
\end{itemize}

This encodes $r_{\text{PM}_{2.5},\text{PM}_{10}}=0.84$ and $r_{\text{NO},\text{NO}_x}=0.90$
by construction rather than asking the model to rediscover them, and it
eliminates an entire class of nonsense output at zero parameter cost.

%======================================================================
\section{Predictive distribution and loss}

\subsection{Why not MSE}

Air quality decisions are threshold decisions: AQI bands, GRAP stages, school
closures. A calibrated upper tail matters more than a good conditional mean, and
an MSE-trained point head systematically under-predicts exactly the
November--December episodes the forecast exists to warn about.

\subsection{Mixture head}

Each query emits, for each of $P$ timesteps, a $K$-component Gaussian mixture in
normalised $\log(1+x)$ space --- equivalently a log-normal mixture in
concentration space:

\begin{equation}
p(y) = \sum_{k=1}^{K}\pi_k\,\mathcal{N}\!\big(y;\ \mu_k,\ \sigma_k^2\big),
\qquad
\pi = \mathrm{softmax}(\ell),\quad \log\sigma_k \in [\sigma_{\min},\sigma_{\max}].
\end{equation}

\begin{equation}
\mathcal{L}_{\text{NLL}} = -\frac{1}{|\Omega|}\sum_{(q,t)\in\Omega}
\log \sum_{k}\exp\!\Big(\log\pi_k -\tfrac{1}{2}\Big(\tfrac{y-\mu_k}{\sigma_k}\Big)^{2} - \log\sigma_k - \tfrac{1}{2}\log 2\pi\Big),
\end{equation}

where $\Omega$ is the set of target positions with an actual observation.

\subsection{CRPS via the order-statistic identity}

NLL alone is unbounded and unstable early in training. We add the continuous
ranked probability score, estimated from $n$ samples $x_1,\dots,x_n$:

\begin{equation}
\mathrm{CRPS} \;\approx\; \frac{1}{n}\sum_i |x_i-y| \;-\; \frac{1}{2n(n-1)}\sum_{i\neq j}|x_i-x_j| .
\end{equation}

The naive pairwise term materialises an $(n,n,\dots)$ tensor and will exhaust
memory at realistic query counts. Using the order statistics
$x_{(1)}\le\cdots\le x_{(n)}$,

\begin{equation}
\boxed{\;\sum_{i,j}|x_i-x_j| \;=\; 2\sum_{i=1}^{n}\big(2i-n-1\big)\,x_{(i)}\;}
\end{equation}

which reduces the cost to $O(n\log n)$ time and $O(n)$ memory. \emph{This is not
an optimisation detail --- the naive form is what makes the loss un-runnable at
scale.}

\subsection{Total objective}

\begin{equation}
\mathcal{L} = w_{\text{nll}}\mathcal{L}_{\text{NLL}}
\;+\; w_{\text{crps}}\mathcal{L}_{\text{CRPS}}
\;+\; w_{\text{phys}}\mathcal{L}_{\text{CRPS}}^{\text{recon}},
\end{equation}

where the third term scores the \emph{reconstructed} observables (PM$_{2.5}$, NO,
NO$_2$) obtained by transforming the sampled ensemble, so the constraint is
scored as a distribution rather than as a point identity. Optionally weight
targets by $1+\alpha\,\sigma\!\big(4(y-y_{\text{thr}})\big)$ to keep winter
episodes from being averaged into irrelevance by clean summer months.

%======================================================================
\section{Pre-training curriculum}

\subsection{Mixed masking}

Random patch masking alone yields a model that is excellent at imputation and
mediocre at forecasting, because there is always future context on both sides of
every gap. The mask distribution is therefore a mixture:

\begin{table}[h]
\centering
\small
\begin{tabular}{@{}llp{7.4cm}@{}}
\toprule
Regime & Weight & Construction \\
\midrule
Tail    & 0.50 & Mask the last $k$ patches, $k\sim\mathcal{U}[1,H]$. No future context: this is forecasting. \\
Block   & 0.35 & Station-wide outages, all channels, durations
$L\sim\mathrm{LogNormal}(\hat\mu,\hat\sigma)$ fitted to the real gap distribution. \\
Channel & 0.15 & Drop a random channel subset over a random window: forces PM$_{2.5}$ to be forecast from meteorology and neighbours alone. \\
\bottomrule
\end{tabular}
\caption{Masking mixture. The block regime is station-wide precisely because the
EDA establishes systemic rather than independent failure.}
\end{table}

Crucially, $\hat\mu,\hat\sigma$ are \emph{estimated from the data}, not guessed:
the observed run-length distribution of zeros in $M$ is exactly the outage
distribution the model must learn to survive.

\subsection{Two-stage schedule}

\begin{enumerate}[leftmargin=1.4em]
\item \textbf{Stage 1 --- national masked pre-training.} All-India corpus, mixed
masking, AdamW ($\beta=(0.9,0.95)$, weight decay $0.05$), cosine schedule with
linear warmup, gradient clipping at $1.0$.
\item \textbf{Stage 2 --- horizon fine-tuning.} Delhi corpus, strict causal split,
tail masking only, fine-tuned at the operational horizons (1\,h, 6\,h, 24\,h,
72\,h) with lower learning rate.
\end{enumerate}

%======================================================================
\section{Reference configuration}

\begin{table}[h]
\centering
\small
\begin{tabular}{@{}llll@{}}
\toprule
Group & Parameter & Value & Note \\
\midrule
Window & context & 512 steps (5.3\,d) & at 15-minute cadence \\
       & patch $P$ & 16 (4\,h) & $N=32$ context patches \\
       & horizon & 288 steps (72\,h) & $H=18$ horizon patches \\
       & neighbourhood & $k=12$ nearest stations & per training sample \\
\midrule
Encoder & $d_{\text{model}}$ & 384 & \\
        & heads & 8 & \\
        & $d_{\text{ff}}$ & 1536 & \\
        & axial blocks & 4 & $=12$ attention layers \\
        & latents/station & 16 & Perceiver bottleneck \\
\midrule
Decoder & layers & 4 & self + cross + FFN \\
        & mixture components $K$ & 3 & log-normal \\
\midrule
Training & optimiser & AdamW, lr $3\times10^{-4}$ & wd $0.05$, clip $1.0$ \\
         & CRPS samples & 32 & order-statistic estimator \\
         & normalisation & masked RevIN & context statistics only \\
\midrule
Size & parameters (Delhi) & $\approx23$M & measured \\
     & parameters (national) & $\approx150$M & $d=768$, 8 blocks \\
\bottomrule
\end{tabular}
\caption{Reference configuration. Activation memory, not FLOPs, is the binding
constraint: the axial encoder stores $\approx1.5$\,GB of activations per sample
at this size, so gradient checkpointing over the encoder blocks (roughly $4\times$
memory reduction for $\sim\!30\%$ step-time cost) is what makes a useful batch
size fit on a single card.}
\end{table}

%======================================================================
\section{Evaluation protocol}

\subsection{Metrics that matter}

Aggregate MAE over a full year is dominated by clean summer months and will look
excellent for a model that is useless in December. Report:

\begin{enumerate}[leftmargin=1.4em]
\item \textbf{Skill stratified by AQI band}, not pooled.
\item \textbf{Episode recall and precision}: detection of
PM$_{2.5}>250\,\mu\text{g}/\text{m}^3$ at 24-hour and 48-hour lead. This is the
number that decides whether the model is operationally useful.
\item \textbf{Calibration}: PIT histograms and reliability of the 90\% and 95\%
predictive intervals, per horizon.
\item \textbf{CRPS by lead time}, to expose where skill decays to climatology.
\item \textbf{Imputation and forecasting scored separately.} They are different
tasks and a single pooled number hides failure in one of them.
\end{enumerate}

\subsection{Baselines}

\begin{table}[h]
\centering
\small
\begin{tabular}{@{}ll@{}}
\toprule
Baseline & Why it is included \\
\midrule
Persistence & Floor; anything below this is broken \\
Seasonal-naive (same hour yesterday) & Captures the diurnal cycle for free \\
Climatology by (station, hour, month) & The honest long-horizon competitor \\
Gradient-boosted trees on lagged features & \textbf{Brutally strong on this task} \\
PatchTST / iTransformer, channel-independent & Isolates the value of spatial attention \\
Design A (per-pollutant decoders) & Isolates the value of the query decoder \\
\bottomrule
\end{tabular}
\caption{If the foundation model does not clearly beat gradient-boosted trees at
24\,h+, the problem is in the setup, not in the baseline.}
\end{table}

\subsection{Splits and leakage discipline}

Hold out 2025--2026 entirely; use rolling-origin backtests within the training
period. Two leakage paths are easy to create and hard to notice:

\begin{itemize}[leftmargin=1.4em]
\item Normalisation statistics computed over the full window rather than the
context. Test this explicitly by perturbing the horizon and asserting the
encoder output is bit-identical.
\item Masked ground truth reaching the model through the mask channel. Assert
that the input mask and target mask never overlap.
\end{itemize}

%======================================================================
\section{Planned ablations}

\begin{table}[h]
\centering
\small
\begin{tabular}{@{}lll@{}}
\toprule
Ablation & Removes & Hypothesis \\
\midrule
No spatial attention & Cross-station information & Large loss during block outages \\
Distance-only bias & Wind asymmetry & Moderate loss at 6--24\,h lead \\
Design A decoders & Query decoder & Coherence violations; no skill gain \\
Random-only masking & Tail regime & Imputation holds, forecasting degrades \\
Point head + MSE & Distributional output & Episode recall collapses \\
No physical reparameterisation & Constraints & PM$_{2.5}>$PM$_{10}$ violations appear \\
Delhi-only pre-training & National corpus & Quantifies the transfer benefit \\
No boundary-layer covariate & ERA5 PBLH & Expected largest single-feature effect \\
\bottomrule
\end{tabular}
\end{table}

%======================================================================
\section{Risks and honest caveats}

\begin{description}[leftmargin=1.4em,style=nextline]
\item[Scale.] $2.16\times10^{8}$ observations is a well-sized deep learning
dataset, not a foundation-model corpus. The ``foundation model'' claim rests on
national pre-training plus demonstrated transfer, and should be stated that way.
\item[Scaling-law transfer.] The $\sim\!20$ tokens/parameter heuristic is
borrowed from language modelling and is not established for time series. Treat
the parameter-count guidance as an order-of-magnitude prior, and run a small
scaling sweep before committing to the national model size.
\item[NWP dependence.] Training with ERA5 as a perfect-prognosis oracle
overstates operational skill. The oracle-versus-forecast gap must be reported.
\item[Non-stationarity.] Emission inventories, GRAP interventions and the vehicle
fleet all change across 2018--2026. A model evaluated only on a held-out tail
may be measuring drift rather than skill; include a rolling re-fit comparison.
\item[Station representativeness.] CPCB station siting is not random; a
kerbside monitor and a background monitor are different measurement processes.
Static covariates $U$ should include a site-type indicator.
\end{description}

%======================================================================
\section{Implementation map}

The accompanying reference implementation is organised as follows.

\begin{table}[h]
\centering
\small
\begin{tabular}{@{}ll@{}}
\toprule
File & Contents \\
\midrule
\texttt{aqfm/config.py}  & \texttt{AQFMConfig}; \texttt{.small()} for tests \\
\texttt{aqfm/layers.py}  & masked RevIN, Fourier time, masked attention, axial block, wind bias \\
\texttt{aqfm/model.py}   & tokeniser, encoder, latent pooler, query decoder, mixture head \\
\texttt{aqfm/losses.py}  & mixture NLL, order-statistic CRPS, episode weighting \\
\texttt{aqfm/masking.py} & mask sampler, empirical outage-length fitting \\
\texttt{aqfm/physical.py}& ratio channels and constrained reconstruction \\
\texttt{aqfm/data.py}    & windowing, neighbourhood sampling, query construction \\
\texttt{aqfm/train.py}   & two-stage training, evaluation, baselines \\
\texttt{tests/}          & shape, leakage, masking, constraint and loss-ordering tests \\
\bottomrule
\end{tabular}
\end{table}

\begin{lstlisting}[caption={Minimal usage.}]
from aqfm import AQFMConfig, AQFM, AirQualityWindows
from aqfm.train import train, evaluate, baselines

cfg = AQFMConfig(grad_checkpoint=True)
ds  = AirQualityWindows(values, observed, timestamps_h, coords, static,
                        cfg, n_neighbours=12, pretrain=True)   # stage 1
model = train(cfg, ds, steps=200_000, batch_size=16, device="cuda")

ds_ft = AirQualityWindows(..., pretrain=False)                 # stage 2
model = train(cfg, ds_ft, model=model, steps=20_000, lr=5e-5)

print(evaluate(model, ds_ft, cfg), baselines(ds_ft, cfg))
\end{lstlisting}

\end{document}
